\chapter{Packet Sniffing and Wireshark}\label{ch:wireshark}
\chapterguide{A basic understanding of hosts, switches, IP traffic, and why packets move through an interface.}{Define packet sniffing and a network traffic trace, navigate the Wireshark interface, capture packets, and connect a packet-list entry to protocol details and raw bytes.}

\section{Purpose and core terms}
Before configuring a network, you need a way to observe its traffic. This lab introduces three linked ideas:
\begin{description}
  \item[Packet sniffing] observing packets as they cross a network.
  \item[Network traffic trace] a time-ordered record of packets seen by an interface.
  \item[Wireshark] a graphical analyzer for inspecting and troubleshooting packet traffic.
\end{description}

\begin{conceptmap}
\textbf{Generate traffic} $\rightarrow$ \textbf{capture on the correct interface} $\rightarrow$ \textbf{find the relevant packet} $\rightarrow$ \textbf{decode protocol fields} $\rightarrow$ \textbf{relate the packet to the action that produced it}.
\end{conceptmap}

\sourceimage[0.64\linewidth]{Packet-sniffing illustration from the supplied slide}{wireshark-sniffing.png}

Capture placement matters: the chosen interface must carry the traffic being investigated.

\section{A traffic trace is evidence, not just a screenshot}
A trace records the sequence of events. If a ping, adjacency, or application session fails, it shows what was sent and whether a reply returned.

\compactsourceimage[0.64\linewidth]{Network-traffic trace concept used in the source material}{wireshark-traffic-trace.png}

\section{The Wireshark workspace}
Wireshark versions differ visually, but the three-pane workflow is consistent.

\sourceimage[0.94\linewidth]{Wireshark packet list, packet details, and packet bytes panes}{wireshark-ui.png}

\begin{tabularx}{\linewidth}{>{\bfseries}p{0.22\linewidth}X}
\toprule
Pane & What you use it for\\\midrule
Packet list & One row per frame, with time, endpoints, protocol, length, and summary.\\
Packet details & Decoded protocol layers, addresses, flags, and fields for the selected frame.\\
Packet bytes & Raw bytes in hexadecimal and ASCII.\\\bottomrule
\end{tabularx}

\section{A disciplined capture workflow}
\begin{netsteps}
\item Identify the traffic you want to observe and the interface that should carry it.
\item Start packet capture on that interface.
\item Generate a small, known event such as a ping so you have a clear reference point.
\item Find the corresponding packet in the packet list using addresses, protocol, and timing.
\item Expand the relevant protocol layers in Packet Details.
\item Consult Packet Bytes only when byte-level evidence is needed.
\item Stop after collecting enough evidence; small traces are easier to analyze.
\end{netsteps}

\begin{workedexample}
If Host A pings Host B, locate the ICMP request and reply. Confirm that the endpoints reverse direction. The trace should prove who sent the request and whether the peer answered.
\end{workedexample}

\section{What to record in a lab submission}
A useful screenshot shows the selected packet, relevant decoded fields, and filter expression. Explain what the evidence proves.

\begin{verifybox}
For this lab, convincing evidence is a capture in which you can point to a known traffic event and identify at least the source, destination, protocol, and the relevant request/reply relationship.
\end{verifybox}

\begin{pitfall}
A capture from the wrong interface, or one buried in unrelated traffic, is poor evidence. Generate controlled traffic and filter by protocol or address.
\end{pitfall}

\begin{quickcheck}
\begin{enumerate}
\item What is the difference between the packet list and packet details panes?
\item Why is a traffic trace easier to interpret when you deliberately generate known traffic?
\item Which pane would you use when you need the exact bytes carried in a frame?
\item If no ping packet appears, which should you question first: the capture location or the packet bytes?
\end{enumerate}
\end{quickcheck}

\begin{chapterrecap}
Wireshark organizes a traffic trace into packet summaries, decoded fields, and raw bytes. This evidence supports verification and troubleshooting throughout the course.
\end{chapterrecap}

\labsection{1}{Capture, filter, and read a live conversation}{sec:lab01-wireshark}

\begin{labproject}{Turn a packet capture into evidence}
\textbf{Coverage.} Chapter~\ref{ch:wireshark}.\\
\textbf{Materials.} \texttt{Lab01/capture\_checklist.md}, \texttt{Lab01/display\_filters.txt}.\\
\textbf{Goal.} Produce a capture that answers a specific question, rather than a screenshot that shows traffic existed.

\begin{labmode}
Capture on a network you are authorised to observe --- your own machine, your own lab segment, or the eNSP simulation. Capturing on a shared campus or corporate network without permission is a disciplinary matter and in many jurisdictions an offence.

Everything below works on loopback or your own interface, so there is no reason to point Wireshark at anyone else's traffic.
\end{labmode}

\medskip\noindent\textbf{Part A --- Ask the question first.}\par
The single most common mistake is starting the capture before deciding what it is meant to prove. You then face 40{,}000 packets and no criterion for which matter.

Write the question down before you press start. Good questions look like:

\begin{itemize}
  \item Did the DNS lookup for \cmd{example.com} succeed, and how long did it take?
  \item Did the TCP handshake complete, or was it reset?
  \item Which host sent the first packet of this conversation?
\end{itemize}

\begin{keyidea}
A capture filter decides what is \emph{recorded}. A display filter decides what is \emph{shown}. They use different syntaxes, and confusing them is the most common source of an empty result window.

Capture filter (BPF): \cmd{tcp port 80}\\
Display filter (Wireshark): \cmd{tcp.port == 80}

Capture broadly and filter afterwards. A capture filter that was too narrow cannot be widened once the traffic has passed.
\end{keyidea}

\medskip\noindent\textbf{Part B --- A disciplined capture.}\par

\begin{netsteps}
  \item Choose the interface carrying the traffic. If unsure, watch the sparklines in the Wireshark start screen.
  \item Start the capture, then generate the traffic --- not the other way round.
  \item Perform exactly one transaction. One page load, one ping sequence, one login.
  \item Stop the capture immediately. Ten seconds of relevant traffic beats ten minutes of noise.
  \item Save the raw \cmd{.pcapng} before filtering. That file is your evidence.
\end{netsteps}

\medskip\noindent\textbf{Part C --- Filters that earn their keep.}\par

\begin{commandmap}
\begin{tabularx}{\linewidth}{@{}>{\ttfamily}p{0.38\linewidth}X@{}}
\toprule
\rowcolor{MLPaleBlue!92}
\textrm{\bfseries Display filter} & \bfseries Shows \\
\midrule
ip.addr == 192.168.1.10 & Everything to or from one host \\
tcp.port == 80 || tcp.port == 443 & Web traffic in both directions \\
tcp.flags.syn == 1 \&\& tcp.flags.ack == 0 & Connection attempts only \\
tcp.analysis.retransmission & Packets Wireshark believes were resent \\
dns & Name resolution \\
icmp & Ping and unreachable messages \\
http.request.method == "GET" & Page requests \\
\bottomrule
\end{tabularx}
\end{commandmap}

\begin{pitfall}
\cmd{ip.addr != 192.168.1.10} does not mean what it appears to. A packet has both a source and a destination address, so the expression is true whenever \emph{either} differs --- which is nearly always.

Write \cmd{!(ip.addr == 192.168.1.10)} instead. The same trap applies to every address and port field.
\end{pitfall}

\medskip\noindent\textbf{Part D --- Read one TCP handshake end to end.}\par
Filter to a single conversation (right-click a packet, then \emph{Follow $\rightarrow$ TCP Stream}) and identify the three-way handshake:

\begin{lstlisting}[style=dcnterminal]
No.  Source        Destination   Proto  Info
1    192.168.1.10  93.184.216.34 TCP    49152 > 80 [SYN] Seq=0 Win=64240
2    93.184.216.34 192.168.1.10  TCP    80 > 49152 [SYN, ACK] Seq=0 Ack=1
3    192.168.1.10  93.184.216.34 TCP    49152 > 80 [ACK] Seq=1 Ack=1
\end{lstlisting}

Answer from your own capture: which host initiated, what ephemeral port did it choose, and how long elapsed between packets 1 and 3? That last figure is the round-trip time, and it is the number a network engineer actually cares about.

\begin{troubleshoot}
\textbf{No packets at all.} Wrong interface, or a capture filter with a typo. Clear the capture filter and try again.

\textbf{Traffic appears but not yours.} You are probably on a switched network seeing only broadcast and your own frames --- which is normal and correct.

\textbf{HTTPS shows no readable content.} That is TLS doing its job. You can still see the handshake, the SNI hostname, and the timing. Analyse those instead.
\end{troubleshoot}

\begin{tryit}
Capture a \cmd{ping} to an address that does not exist on your subnet. No reply arrives --- so what \emph{does} the capture contain, and what does that tell you about how the host tried to find the destination? Look for ARP.
\end{tryit}
\end{labproject}

\begin{verifybox}
\begin{itemize}
  \item The raw \cmd{.pcapng} is saved and submitted, not only screenshots.
  \item The question the capture answers is stated before the evidence.
  \item Display filters are quoted exactly as used.
  \item Packet numbers are cited when referring to specific packets.
  \item The capture was taken on a network you are authorised to observe.
\end{itemize}
\end{verifybox}

\begin{labdeliverable}
Submit the \cmd{.pcapng} file, the question it answers, the display filter used, an annotated screenshot of the three-way handshake with packet numbers, and the measured round-trip time. Add one sentence on why capturing broadly and filtering afterwards is the safer order.
\end{labdeliverable}

